\documentclass[11pt]{article}
\usepackage[margin=1in]{geometry}
\usepackage{amsmath,amssymb}
\usepackage{enumitem}
\usepackage[T1]{fontenc}

\title{COMP2300 Set 6 Practice Exam}
\author{Topics: Branch-and-Link, Calling Conventions, Stack, ABI, Recursion}
\date{}

\begin{document}
\maketitle

\noindent\textbf{Time:} 120 minutes \hfill
\textbf{Total:} 100 marks

\vspace{0.5em}
\noindent\textbf{Instructions}
\begin{itemize}[leftmargin=1.5em]
  \item Answer in English.
  \item Part A is concept-focused multiple choice.
  \item Part B contains computational and trace-based questions.
  \item Part C contains adapted past-exam style questions with sources.
  \item Assume 32-bit ARM unless stated otherwise.
\end{itemize}

\section*{Part A: Multiple Choice (30 marks, 3 marks each)}
Choose exactly one option for each question.

\begin{enumerate}[label=\textbf{A\arabic*.}]
  \item The main extra job performed by \texttt{BL} compared with a regular branch is to:
  \begin{enumerate}[label=(\Alph*)]
    \item Zero the stack pointer
    \item Save the return address in \texttt{LR}
    \item Save all caller registers automatically
    \item Write the target into \texttt{R0}
  \end{enumerate}

  \item In the course convention, the first four function arguments are typically passed in:
  \begin{enumerate}[label=(\Alph*)]
    \item \texttt{R4-R7}
    \item \texttt{R8-R11}
    \item \texttt{R0-R3}
    \item \texttt{R12-R15}
  \end{enumerate}

  \item Which register conventionally holds a function return value?
  \begin{enumerate}[label=(\Alph*)]
    \item \texttt{R0}
    \item \texttt{R4}
    \item \texttt{R12}
    \item \texttt{LR}
  \end{enumerate}

  \item Which statement about the stack in the lecture material is correct?
  \begin{enumerate}[label=(\Alph*)]
    \item It is a FIFO queue
    \item It is used only for recursion
    \item It is used to save registers, pass extra arguments, and hold local data
    \item It is part of the ISA but not the ABI
  \end{enumerate}

  \item Which pair best describes the division of responsibility in a calling convention?
  \begin{enumerate}[label=(\Alph*)]
    \item Hardware-save and software-save
    \item Caller-save and callee-save
    \item Decode-save and execute-save
    \item Stack-save and heap-save
  \end{enumerate}

  \item Which register is the stack pointer?
  \begin{enumerate}[label=(\Alph*)]
    \item \texttt{R11}
    \item \texttt{R12}
    \item \texttt{R13}
    \item \texttt{R14}
  \end{enumerate}

  \item Which statement about \texttt{ABI} is correct?
  \begin{enumerate}[label=(\Alph*)]
    \item It is identical to the ISA
    \item It includes procedure-call rules such as calling convention
    \item It is only about cache organization
    \item It defines transistor-level behavior
  \end{enumerate}

  \item A recursive function is always:
  \begin{enumerate}[label=(\Alph*)]
    \item A leaf function
    \item A function with no local variables
    \item A non-leaf function that calls itself
    \item A function that does not use the stack
  \end{enumerate}

  \item Which of the following is preserved across a correct function call under a calling convention?
  \begin{enumerate}[label=(\Alph*)]
    \item Every temporary register automatically
    \item Every argument register automatically
    \item Any callee-saved register that the callee uses, after it is restored
    \item The entire data memory
  \end{enumerate}

  \item A stack frame is:
  \begin{enumerate}[label=(\Alph*)]
    \item The set of all registers in the CPU
    \item The portion of stack space allocated for one active function call
    \item The instruction encoding of \texttt{BL}
    \item The branch target address
  \end{enumerate}
\end{enumerate}

\section*{Part B: Computational Problems (50 marks)}
\begin{enumerate}[label=\textbf{B\arabic*.}]
  \item \textbf{Branch target address (8 marks)}\\
  An instruction at address \texttt{0x2040} is a branch with \texttt{imm24 = 0x000005}.
  \begin{enumerate}[label=(\alph*)]
    \item Compute \texttt{imm24 << 2}.
    \item Compute the sign-extended offset.
    \item Compute the branch target address using the ARM rule from the lecture.
  \end{enumerate}

  \item \textbf{Register roles in a call (8 marks)}\\
  Consider the function call:
\begin{verbatim}
MOV R0, #10
MOV R1, #20
BL foo
\end{verbatim}
  The caller still needs the old values of \texttt{R0} and \texttt{R1} after the call.
  \begin{enumerate}[label=(\alph*)]
    \item Which register holds the return address after \texttt{BL foo}?
    \item Are \texttt{R0} and \texttt{R1} caller-saved or callee-saved in the lecture convention?
    \item What must the caller do before the call if it still needs those values afterwards?
  \end{enumerate}

  \item \textbf{Stack-pointer updates (8 marks)}\\
  Assume the stack is full descending and initially \texttt{SP = 0x1000}.
  \begin{enumerate}[label=(\alph*)]
    \item What is \texttt{SP} after \texttt{PUSH \{R4, R5, LR\}}?
    \item What is \texttt{SP} after then executing \texttt{SUB SP, SP, \#8} for two 32-bit locals?
    \item If the function later executes \texttt{ADD SP, SP, \#8} and \texttt{POP \{R4, R5, LR\}}, what is the final \texttt{SP}?
  \end{enumerate}

  \item \textbf{Prologue and epilogue reasoning (8 marks)}\\
  A function uses registers \texttt{R4} and \texttt{R7}, returns in \texttt{R0}, and makes a nested function call.
  \begin{enumerate}[label=(\alph*)]
    \item Explain why \texttt{LR} must be preserved.
    \item Which of \texttt{R4} and \texttt{R7} must be saved if the function modifies them?
    \item Write a possible prologue and epilogue using \texttt{PUSH}/\texttt{POP}.
  \end{enumerate}

  \item \textbf{Stack frames in recursion (8 marks)}\\
  A recursive function saves \texttt{R0} and \texttt{LR} using \texttt{PUSH \{R0, LR\}} on each call.
  Assume the initial \texttt{SP = 0x8000} and the function is called with \texttt{n = 3}, leading to active calls for \texttt{n=3}, \texttt{n=2}, and \texttt{n=1}.
  \begin{enumerate}[label=(\alph*)]
    \item How many active stack frames exist at the deepest point?
    \item How many bytes have been consumed by those saves alone?
    \item What is \texttt{SP} at the deepest point?
  \end{enumerate}

  \item \textbf{Caller-save vs callee-save (10 marks)}\\
  Suppose function \texttt{main} stores a useful temporary in \texttt{R2}, calls \texttt{f}, and later still needs the old value of \texttt{R2}. Function \texttt{f} modifies \texttt{R4}, \texttt{R5}, and \texttt{R2}, and then calls \texttt{g}.
  \begin{enumerate}[label=(\alph*)]
    \item Which function is responsible for preserving the old value of \texttt{R2} across the call to \texttt{f}?
    \item Which function is responsible for preserving \texttt{R4} and \texttt{R5} if \texttt{f} changes them?
    \item Why is \texttt{f} also responsible for preserving \texttt{LR} before calling \texttt{g}?
    \item Briefly explain why this division of responsibility is more efficient than ``save every register every time''.
  \end{enumerate}
\end{enumerate}

\section*{Part C: Past Exam Questions (Source-Tagged) (20 marks)}
\begin{enumerate}[label=\textbf{C\arabic*.}]
  \item \textbf{Functions with stack support (Source: 2023\_Final\_Exam\_SuppDA.pdf, Part A, adapted) (6 marks)}\\
  A teaching ISA is extended with \texttt{sp}, \texttt{lr}, \texttt{push}, \texttt{pop}, \texttt{bl}, and \texttt{bx lr}. Explain why each of the following is useful for implementing functions:
  \begin{enumerate}[label=(\alph*)]
    \item An explicit link register
    \item Stack push/pop support
    \item A branch-and-link style call instruction
  \end{enumerate}

  \item \textbf{Calling convention design (Source: 2023\_Final\_Exam\_SuppDA.pdf, Part A, adapted) (7 marks)}\\
  Suppose a simple ISA has a function \texttt{mul} that may clobber registers. Propose a small calling convention for four general-purpose registers and state clearly which are caller-saved and which are callee-saved. Briefly justify your design.

  \item \textbf{Returning without a link register (Source: 2023\_Final\_Exam.pdf, Q1A, adapted) (7 marks)}\\
  Explain how functions can still be implemented in an ISA that has a stack pointer but no dedicated link register. What must be saved on the stack at call time, and how is the return performed?
\end{enumerate}

\end{document}
